\chapter{Elementary functions from power series}
This chapter uses the general series results of modules 11--13 as prior logical
ingredients, irrespective of the printed chapter order. Those general results
use factorials, finite algebra and completeness, and do not use the exponential
or logarithm constructed here. Throughout, derivatives of complex-valued
functions of a real parameter use real increments.
\SourceRef{Rudin, Chapter 8, PDF pp. 183--191; Zorich I, Sections 5.5 and 5.6.5. The series-first route is fixed here; integral and differential-equation descriptions are comparisons}

\KEntry{K03.07.exp_construction}{The exponential series}
\begin{construction}
For $z\in\C$ define
\[E(z)=\sum_{n=0}^{\infty}\frac{z^n}{n!},\qquad 0!=1.
\]
Its restriction to $\R$ will also be written $\exp$.
\end{construction}
\begin{proof}
For $|z|\le R$, the absolute values are bounded by $R^n/n!$.
When $n+1\ge2R$, successive terms of this majorant have ratio at most
$1/2$, so its tail is bounded by a geometric tail. Thus the series converges
absolutely everywhere and uniformly on every disk $|z|\le R$ by the
uniform majorant argument of \KRef{K03.12.m_test}. The differentiated
series on a real interval is $\sum_{n\ge1}x^{n-1}/(n-1)!$ and obeys the
same bound. The general termwise theorem \KRef{K03.13.termwise} therefore
gives $E'(x)=E(x)$. This reasoning is about a previously established general
power-series theorem, not a definition of differentiation by formal symbols.
\end{proof}

\KEntry{K03.07.exp_addition}{Addition law}
\begin{theorem}
For $z,w\in\C$, $E(z+w)=E(z)E(w)$, $E(0)=1$, and
$E(-z)=1/E(z)$.
\end{theorem}
\begin{proof}
The absolute Cauchy-product theorem \KRef{K03.11.cauchy_product} applies
because both factorial series converge absolutely. Its coefficient of total
degree $n$ is
\[\sum_{k=0}^n\frac{z^kw^{n-k}}{k!(n-k)!}
=\frac1{n!}\sum_{k=0}^n\binom nkz^kw^{n-k}
=\frac{(z+w)^n}{n!}.
\]
The constant value is read from the series. Setting $w=-z$ gives
$E(z)E(-z)=1$, proving both nonvanishing and the reciprocal formula.
\end{proof}

\KEntry{K03.07.exp_properties}{The real exponential}
\begin{theorem}
The real function $\exp$ is positive, differentiable with derivative itself,
strictly increasing, and maps $\R$ bijectively onto $(0,\infty)$. Its limits
at $-\infty$ and $+\infty$ are $0$ and $+\infty$, respectively.
\end{theorem}
\begin{proof}
For $x\ge0$ the series has nonnegative terms and $E(x)\ge1+x>0$.
For $x<0$, $E(x)=1/E(-x)>0$. The derivative proved above is consequently
positive, so the mean value theorem makes $E$ strictly increasing.
The bound $E(x)\ge1+x$ proves divergence to $+\infty$; the reciprocal law
proves the limit at $-\infty$. Continuity and the intermediate value theorem
then show that each positive number is assumed exactly once.
\end{proof}

\KEntry{K03.07.log_construction}{The inverse logarithm}
\begin{definition}
For $y>0$, $\log y$ is the unique real $x$ such that $E(x)=y$.
\end{definition}
\begin{proposition}
The logarithm is continuous and strictly increasing on $(0,\infty)$,
$\log1=0$, and $\log(uv)=\log u+\log v$ for $u,v>0$.
\end{proposition}
\begin{proof}
The inverse-continuity theorem \KRef{K03.03.inverse_continuity} applies to
the continuous strictly increasing surjection just proved. Alternatively,
if $x=\log y$, the images of $x-\eps$ and $x+\eps$ enclose $y$ and trap
the inverse of every sufficiently close positive number between these two
points. Finally $E(\log u+\log v)=uv$ by addition, and inverse uniqueness
proves the product formula. In particular $\log(1/u)=-\log u$.
\end{proof}

\KEntry{K03.07.log_derivative}{Derivative and integral comparison for logarithm}
\begin{theorem}
For $x>0$, $\log'(x)=1/x$. After the fundamental theorem in module 09,
\[\log x=\int_1^x\frac{\dd t}{t}\quad(x>0)
\]
with oriented integration.
\end{theorem}
\begin{proof}
The inverse derivative rule applies because $E'(\log x)=E(\log x)=x\ne0$;
it gives the asserted reciprocal derivative. On every compact interval in
$(0,\infty)$, $1/t$ is continuous, so \KRef{K03.09.ftc_first} makes the
accumulation integral a primitive. Its difference from $\log$ has derivative
zero, hence is constant on $(0,\infty)$. Both functions vanish at $1$.
Only this comparison uses integration; the construction of $\log$ does not.
\end{proof}

\KEntry{K03.07.real_powers}{Real powers and roots}
\begin{definition}
For $x>0$ and $\alpha\in\R$, put $x^\alpha=E(\alpha\log x)$.
\end{definition}
\begin{proposition}
For positive bases,
\[x^{\alpha+\beta}=x^\alpha x^\beta,\quad
(xy)^\alpha=x^\alpha y^\alpha,\quad (x^\alpha)^\beta=x^{\alpha\beta}.
\]
This agrees with integer powers and with the positive algebraic root
interpretation of rational powers.
\end{proposition}
\begin{proof}
The first two laws follow from exponential addition and logarithm products;
the third follows from $\log(E(\alpha\log x))=\alpha\log x$.
For integers the first law, $x^1=x$, and $x^0=1$ give the familiar positive
and negative powers. If $q\ge1$ and $p\in\Z$, then
$(x^{p/q})^q=x^p$ and $x^{p/q}>0$, so uniqueness of the positive $q$th
root from \KRef{K03.01.positive_roots} proves agreement.
For $\alpha>0$ the limit at $x\downarrow0$ is $0$, since $\log x\to-\infty$;
thus $0^\alpha=0$ is a continuous extension. For $\alpha<0$ the limit is
$+\infty$. No definition of $0^0$ follows. Negative bases admit integer
powers, and rational powers with odd reduced denominator by real roots;
there is no real-valued extension for all real exponents at a negative base.
\end{proof}

\KEntry{K03.07.sin_cos}{Sine and cosine series}
\begin{construction}
For real or complex $z$ define
\[S(z)=\sum_{n=0}^{\infty}\frac{(-1)^nz^{2n+1}}{(2n+1)!},\qquad
C(z)=\sum_{n=0}^{\infty}\frac{(-1)^nz^{2n}}{(2n)!}.
\]
Write $\sin=S$ and $\cos=C$. On real inputs $S'=C$ and $C'=-S$.
\end{construction}
\begin{proof}
Absolute values of either series are subseries of $\sum |z|^n/n!$;
on $|z|\le R$ the same majorant proves uniform convergence. Their
differentiated real series are the displayed cosine and negative sine
series, with the same compact majorants. Termwise differentiation yields
the formulas, with $S(0)=0$ and $C(0)=1$.
\end{proof}

\KEntry{K03.07.trig_properties}{Trigonometric identities and bounds}
\begin{theorem}
Sine is odd, cosine is even, and for real $x,y$,
\begin{align*}
S(x+y)&=S(x)C(y)+C(x)S(y),\\
C(x+y)&=C(x)C(y)-S(x)S(y),\\
S(x)^2+C(x)^2&=1.
\end{align*}
Consequently $|S(x)|,|C(x)|\le1$, $|S(x)|\le|x|$, and
$|C(x)-1|\le x^2/2$.
\end{theorem}
\begin{proof}
Parity follows term by term. To prove the addition formulas without assuming
geometric trigonometry, fix $y$ and subtract the proposed expressions from
$S(x+y)$ and $C(x+y)$, calling the differences $u(x),v(x)$.
The derivative formulas give $u'=v$, $v'=-u$, and $u(0)=v(0)=0$.
Hence $(u^2+v^2)'=0$, and its initial value zero implies $u=v=0$.
Likewise $(S^2+C^2)'=2SC-2CS=0$, with value one at zero.
The norm bounds follow, and the mean value inequality gives $|S(x)|\le|x|$.
For the sharper cosine estimate, put $h(x)=x^2/2-1+C(x)$ for $x\ge0$.
Then $h'=x-S(x)\ge0$, so $1-C(x)\le x^2/2$; evenness handles $x<0$.
\end{proof}

\KEntry{K03.07.pi_period}{Pi, periods and quadrants}
\begin{construction}
Let $c$ be the first positive zero of $C$ and define $\pi=2c$.
Then $S(c)=1$, $C(c)=0$, sine and cosine have period $2\pi$,
and their first-quadrant ranges are $S([0,c])=[0,1]$ and
$C([0,c])=[0,1]$ with strict increase and decrease, respectively.
\end{construction}
\begin{proof}
Continuity makes $C>0$ on some $[0,\delta]$, and then $S$ increases there,
with $S(\delta)>0$. If $C$ had no positive zero, continuity would force
$C>0$ everywhere on $(0,\infty)$. Thus $S(x)\ge S(\delta)>0$ for
$x\ge\delta$ and $C'(x)=-S(x)\le-S(\delta)$; the mean value theorem
would make $C$ negative for large $x$, a contradiction. A zero exists.
The zero set in a compact interval from $\delta$ to any such zero is
nonempty and closed, hence has a minimum $c>0$. On $(0,c)$, $C>0$,
so $S$ is positive and strictly increasing. The identity gives $S(c)=1$.
Since $C'=-S<0$ on $(0,c]$, the range statements follow.
Addition gives $S(x+c)=C(x)$, $C(x+c)=-S(x)$; applying these twice gives
$S(x+2c)=-S(x)$, $C(x+2c)=-C(x)$ and four times gives period $4c=2\pi$.
These relations and parity describe all quadrants: sine increases on
$[-c,c]$, decreases on $[c,3c]$, and cosine decreases on $[0,2c]$,
with translated intervals giving all other monotonicity intervals.
They also show that the zeros of sine are $k\pi$ and those of cosine
$c+k\pi$ ($k\in\Z$). There is no common positive period below $2\pi$:
a common period must have sine zero and cosine one, which these quadrant
signs allow only at even multiples of $\pi$.
This analytic definition does not assume a circle-length definition of $\pi$.
\end{proof}

\KEntry{K03.07.trig_inverses}{Principal inverse trigonometric functions}
\begin{definition}
$\arcsin:[-1,1]\to[-\pi/2,\pi/2]$ is the inverse of sine on that range;
$\arccos:[-1,1]\to[0,\pi]$ is the inverse of cosine there.
$\arctan:\R\to(-\pi/2,\pi/2)$ is the inverse of tangent on that interval.
\end{definition}
\begin{proof}[Existence and derivative formulas]
The preceding strict monotonicity and endpoint values prove the first two
bijections and inverse continuity. On $(-\pi/2,\pi/2)$ cosine is positive;
the quotient rule gives $(S/C)'=1/C^2>0$. At the endpoints its limits are
$-\infty,+\infty$ because the numerator tends to $-1,1$ and the
denominator to $0$ through positive values. Thus tangent is a continuous
strictly increasing surjection, establishing the third inverse.
For $|x|<1$, inverse differentiation and the sign of the appropriate sine
or cosine give
\[(\arcsin x)'=\frac1{\sqrt{1-x^2}},\qquad
(\arccos x)'=-\frac1{\sqrt{1-x^2}}.
\]
For all $x\in\R$, $(\arctan x)'=1/(1+x^2)$ because
$C(\arctan x)^{-2}=1+x^2$. At $x=\pm1$ the first two inverse functions
remain continuous, but their finite one-sided derivatives do not exist:
the interior derivative magnitudes tend to infinity and the mean value
theorem on a one-sided secant gives the same divergence of slopes.
\end{proof}

\KEntry{K03.07.tan}{Tangent and poles}
\begin{proposition}
$\tan x=S(x)/C(x)$ is defined exactly for
$x\notin\{\pi/2+k\pi:k\in\Z\}$, has period $\pi$, and derivative
$1/C(x)^2$. At each pole $p=\pi/2+k\pi$, the left and right limits
are $+\infty$ and $-\infty$.
\end{proposition}
\begin{proof}
The zero description for $C$ proves the domain. Both numerator and
denominator change sign after $\pi$, proving periodicity.
The quotient rule gives $(C^2+S^2)/C^2=1/C^2$.
At $\pi/2$, sine tends to $1$ while cosine is positive on the left and
negative on the right. Periodicity translates these one-sided limits to
every pole. A pole is excluded from the domain; the limits are not values.
\end{proof}

\KEntry{K03.07.hyperbolic}{Hyperbolic functions and their real inverses}
\begin{definition}
Put $\sinh x=(E(x)-E(-x))/2$, $\cosh x=(E(x)+E(-x))/2$ and
$\tanh x=\sinh x/\cosh x$ for $x\in\R$.
\end{definition}
\begin{proposition}
$\cosh^2x-\sinh^2x=1$, $(\sinh)'=\cosh$, $(\cosh)'=\sinh$,
and $(\tanh)'=1/\cosh^2$. The respective principal inverses are
\begin{align*}
\operatorname{arsinh}y&=\log(y+\sqrt{y^2+1})&& (y\in\R),\\
\operatorname{arcosh}y&=\log(y+\sqrt{y^2-1})&& (y\ge1),\\
\operatorname{artanh}y&=\tfrac12\log\frac{1+y}{1-y}&& (|y|<1),
\end{align*}
with ranges $\R,[0,\infty),\R$. Their derivatives on the interiors
of these domains are $(1+y^2)^{-1/2}$, $(y^2-1)^{-1/2}$,
and $(1-y^2)^{-1}$.
\end{proposition}
\begin{proof}
Expanding the squares and using $E(x)E(-x)=1$ proves the identity;
differentiation proves the three derivative formulas. Since $\cosh x>0$,
$\sinh$ increases strictly and tends to both real infinities. Cosine
hyperbolic decreases to $1$ on $(-\infty,0]$ and increases from $1$ on
$[0,\infty)$. The quotient expression
$\tanh x=(E(2x)-1)/(E(2x)+1)$ gives strict increase and limits $\pm1$.
Solving each equation using $u=E(x)>0$ gives the displayed logarithms;
for example $2y=u-u^{-1}$ has the unique positive root
$u=y+\sqrt{y^2+1}$. For cosh select $u\ge1$ to fix the nonnegative branch.
Inverse differentiation gives the stated derivatives from the identity,
with the positive square root on the chosen branch. At $y=1$ arcosh
is continuous with infinite right derivative; no finite derivative is claimed.
\end{proof}

\KEntry{K03.07.euler}{Imaginary inputs and Euler's identity}
\begin{theorem}
For real $t$, $E(it)=C(t)+iS(t)$, $|E(it)|=1$, and
$\overline{E(z)}=E(\bar z)$ for every $z\in\C$. For fixed $z\in\C$,
the real-parameter derivative of $E(tz)$ is $zE(tz)$.
\end{theorem}
\begin{proof}
Absolute convergence permits separating the even and odd terms in $E(it)$,
which gives the identity. The squared modulus is $C^2+S^2=1$.
Conjugation commutes with finite partial sums and limits, giving the second
identity. Uniform convergence of the differentiated series on bounded real
$t$-intervals gives
$\sum_{n\ge1}nz^nt^{n-1}/n!=zE(tz)$ componentwise. This asserts a real
parameter derivative; it does not invoke complex differentiability theory.
\end{proof}

\KEntry{K03.07.elementary_derivatives}{A derivative table with domains}
\begin{proposition}
For $x>0$ and $\alpha\in\R$,
\[(x^\alpha)'=\alpha x^{\alpha-1},\qquad
(a^x)'=(\log a)a^x\quad(a>0).
\]
Together with the formulas proved above this gives the derivatives of the
elementary functions on their stated domains. In particular, $\sec=1/\cos$
has derivative $\sec\tan$ where $\cos\ne0$, and $\cot=\cos/\sin$ has
derivative $-1/\sin^2$ where $\sin\ne0$.
\end{proposition}
\begin{proof}
The chain rule gives $(E(\alpha\log x))'=E(\alpha\log x)\alpha/x$,
equal to $\alpha x^{\alpha-1}$ by the power laws. The second formula is
the chain rule for $E(x\log a)$. The reciprocal and quotient rules give
the last two formulas. Thus every entry follows from an actual derivative
calculation with its denominator condition, rather than a formal table
applied across zeros or branch endpoints.
\end{proof}

\KEntry{K03.07.elementary_domains}{Domains and exceptional points}
\begin{definition}
In these notes $\log$ has real domain $(0,\infty)$, $x^\alpha$ initially
has $x>0$, $\sqrt{x}$ is the nonnegative root for $x\ge0$, and the three
inverse trigonometric ranges are exactly those fixed above. Reciprocal
trigonometric functions exclude their denominator zeros.
\end{definition}
\begin{example}
The equation $\sin y=x$ has a unique solution in $[-\pi/2,\pi/2]$, but
not a unique real solution without that restriction: for $|x|<1$ both
$\arcsin x$ and $\pi-\arcsin x$ solve it, along with their $2\pi$
translates. Also $(\sqrt{x})'=1/(2\sqrt{x})$ holds for $x>0$;
the right quotient at zero is $1/\sqrt{h}\to+\infty$.
These computations explain the exclusions. A complex logarithm and its
branches require a different domain theory and are outside this chapter.
\end{example}

\KEntry{K03.07.standard_limits}{Reference limits}
\begin{theorem}
As $h\to0$ through real nonzero values,
\[\frac{\sin h}{h}\to1,\qquad\frac{E(h)-1}{h}\to1,\qquad
\frac{\log(1+h)}h\to1\quad(1+h>0).
\]
More generally, if $u\to0$, $1+u>0$, and $v u\to L\in\R$, then
$(1+u)^v\to E(L)$, provided at zeros $u=0$ the expression is interpreted
as $1$ and the asserted product limit is consistent there.
\end{theorem}
\begin{proof}
The three quotients are derivative quotients at $0,0,1$, respectively,
whose derivatives have been proved to be one. For the last assertion,
at nonzero $u$,
$v\log(1+u)=(vu)\log(1+u)/u\to L$.
At zeros the exponent is zero; if zeros accumulate, existence of the
product limit forces $L=0$ along that subsequence. Continuity of $E$
completes the proof, including this case.
In particular $\lim_{h\to0}(1+h)^{1/h}=E(1)$ on the positive-base
domain, by taking $u=h$ and $v=1/h$ for nonzero $h$.
\end{proof}

\KEntry{K03.07.construction_bridges}{Sequence, differential equation and integral comparisons}
\begin{theorem}
For every fixed $x\in\R$,
\[\lim_{n\to\infty}(1+x/n)^n=E(x).
\]
On an interval containing zero, $E$ is the unique differentiable scalar
solution of $y'=y$, $y(0)=1$. The integral logarithm agrees with the
inverse logarithm.
\end{theorem}
\begin{proof}
For all sufficiently large $n$, $1+x/n>0$. If $x\ne0$, then
$n\log(1+x/n)=x\,\log(1+x/n)/(x/n)\to x$ by the preceding limit;
exponentiating proves the sequence assertion. For $x=0$ every term is one.
For the equation, existence is $E'=E$, $E(0)=1$. If $y$ is another
real or complex solution, the real derivative of $y(t)E(-t)$ is
$y'(t)E(-t)-y(t)E(-t)=0$. The zero-derivative theorem, componentwise
for complex values, gives $y(t)E(-t)=1$ on the interval.
Finally the integral comparison is precisely the primitive-uniqueness
argument of \KRef{K03.07.log_derivative}. These are proved identifications,
not simultaneous definitions that would create a circular construction.
\end{proof}
